\documentclass[../../main.tex]{subfiles}% 注意这里的文件路径不能用 ./main.tex ,否则用latexmk编译子文件会报错
\graphicspath{{\subfix{./image/}}{\subfix{../../image/}}} % 指定图片引用路径,后续可以直接使用图片文件名
% 如果图片存放在上级目录,则使用 ../../image/ 作为备用路径

% 例如:
% \begin{figure}[H]
% \centering
% \includegraphics[width=0.3\textwidth]{图.png}
% \caption{}
% \label{figure:图}
% \end{figure}
% 注意:上述\label{}一定要放在\caption{}之后,否则引用图片序号会只会显示??.

\begin{document}

\section{乘积测度}\label{section:ch02-乘积测度}

\subsection{Tonelli 与 Fubini 定理}\label{subsection:ch02-TonelliFubini}

设 $(X_1,\mathscr A_1,\mu_1)$ 与 $(X_2,\mathscr A_2,\mu_2)$ 为测度空间。设 $E\in\mathscr A_1\times\mathscr A_2$。于是对任意 $x_1\in X_1$ 有 $E_{x_1}\in\mathscr A_2$，对任意 $x_2\in X_2$ 有 $E_{x_2}\in\mathscr A_1$（见\cref{proposition:ch01-截面}）。则从 $X_1$ 到 $[0,\infty]$ 的映射 $x_1\mapsto\mu_2(E_{x_1})$（$x_1\in X_1$），以及从 $X_2$ 到 $[0,\infty]$ 的映射 $x_2\mapsto\mu_1(E_{x_2})$（$x_2\in X_2$）均良定义。有下列结果。

\begin{proposition}\label{proposition:ch02-乘积截面积分}
设 $(X_1,\mathscr A_1,\mu_1)$ 与 $(X_2,\mathscr A_2,\mu_2)$ 为 $\sigma$-有限测度空间。设 $E\in\mathscr A_1\times\mathscr A_2$。则：
\begin{enumerate}[(i)]
\item 映射 $X_1\to[0,\infty]$，$x_1\mapsto\mu_2(E_{x_1})$，为 $\mathscr A_1$-可测；映射 $X_2\to[0,\infty]$，$x_2\mapsto\mu_1(E_{x_2})$，为 $\mathscr A_2$-可测；
\item 有
\begin{align}
\int_{X_1}\mu_2(E_{x_1})\,\mathrm{d}\mu_1=\int_{X_2}\mu_1(E_{x_2})\,\mathrm{d}\mu_2; \label{align:ch02-2.15a}
\end{align}
\item 由
\begin{align}
(\mu_1\times\mu_2)(E)\triangleq \int_{X_1}\mu_2(E_{x_1})\,\mathrm{d}\mu_1=\int_{X_2}\mu_1(E_{x_2})\,\mathrm{d}\mu_2\quad(E\in\mathscr A_1\times\mathscr A_2) \label{align:ch02-2.15b}
\end{align}
定义的映射 $\mu_1\times\mu_2:\mathscr A_1\times\mathscr A_2\to[0,\infty]$ 为 $\sigma$-有限测度。
\end{enumerate}
\end{proposition}

\begin{remark}\label{remark:ch02-sigma有限必要}
若 $\mu_1$ 与/或 $\mu_2$ 非 $\sigma$-有限，易构造使\eqref{align:ch02-2.15a}不成立的例子。
\end{remark}

由\cref{proposition:ch02-乘积截面积分}，下列定义是适定的。

\begin{definition}\label{definition:ch02-乘积测度}
设 $(X_1,\mathscr A_1,\mu_1)$ 与 $(X_2,\mathscr A_2,\mu_2)$ 为 $\sigma$-有限测度空间。由\eqref{align:ch02-2.15b}定义的测度 $\mu_1\times\mu_2$ 称为乘积测度，$(X_1\times X_2,\mathscr A_1\times\mathscr A_2,\mu_1\times\mu_2)$ 称为乘积测度空间。
\end{definition}

\begin{remark}\label{remark:ch02-乘积a.e.等价}
设 $(X_1\times X_2,\mathscr A_1\times\mathscr A_2,\mu_1\times\mu_2)$ 为乘积测度空间，$P=P(x_1,x_2)$ 为某性质。利用\eqref{align:ch02-2.15b}易见下列命题等价：(i) $P$ 在 $X_1\times X_2$ 中 $\mu_1\times\mu_2$-a.e. 成立；(ii) 对 $\mu_1$-a.e. $x_1\in X_1$，$P(x_1,\cdot)$ 在 $X_2$ 中 $\mu_2$-a.e. 成立；(iii) 对 $\mu_2$-a.e. $x_2\in X_2$，$P(\cdot,x_2)$ 在 $X_1$ 中 $\mu_1$-a.e. 成立。
\end{remark}

\begin{remark}\label{remark:ch02-乘积完备化}\label{proposition:ch02-完备化}
简单例子表明完备测度空间的乘积不必完备。然而可证 Lebesgue 空间 $(\mathbb R^{M+N},\mathscr L_{M+N},\lambda_{M+N})$ 是乘积 $(\mathbb R^{M+N},\mathscr L_M\times\mathscr L_N,\lambda_M\times\lambda_N)$ 的完备化（$M,N\in\mathbb N$；见\cref{definition:ch01-完备化}）。
\end{remark}

设 $(X_1,\mathscr A_1,\mu_1)$ 与 $(X_2,\mathscr A_2,\mu_2)$ 为测度空间，$f:X_1\times X_2\to[0,\infty)$ 为 $\mathscr A_1\times\mathscr A_2$-可测。由\cref{proposition:ch02-迹可测}，$x_1$-迹 $f_{x_1}$ 为 $\mathscr A_2$-可测，故函数
\begin{align}
X_1\to[0,\infty],\qquad x_1\mapsto\int_{X_2}f_{x_1}\,\mathrm{d}\mu_2 \label{align:ch02-2.16a}
\end{align}
良定义；类似地，映射
\begin{align}
X_2\to[0,\infty],\qquad x_2\mapsto\int_{X_1}f_{x_2}\,\mathrm{d}\mu_1 \label{align:ch02-2.16b}
\end{align}
良定义。

\begin{theorem}[Tonelli]\label{theorem:ch02-Tonelli}
设 $(X_1,\mathscr A_1,\mu_1)$ 与 $(X_2,\mathscr A_2,\mu_2)$ 为 $\sigma$-有限测度空间。设 $f:X_1\times X_2\to[0,\infty)$ 为 $\mathscr A_1\times\mathscr A_2$-可测。则：
\begin{enumerate}[(i)]
\item \eqref{align:ch02-2.16a}中定义的函数分别为 $\mathscr A_1$-、$\mathscr A_2$-可测；
\item 有
\begin{align}
\int_{X_1\times X_2}f\,\mathrm{d}(\mu_1\times\mu_2)
=\int_{X_1}\mathrm{d}\mu_1\int_{X_2}f_{x_1}\,\mathrm{d}\mu_2
=\int_{X_2}\mathrm{d}\mu_2\int_{X_1}f_{x_2}\,\mathrm{d}\mu_1. \label{align:ch02-2.17}
\end{align}
\end{enumerate}
\end{theorem}

\begin{remark}\label{remark:ch02-Tonelli条件}
\label{remark:ch02-截口注}
注意\eqref{align:ch02-2.15b}是\eqref{align:ch02-2.17}在 $f=\chi_E$（$E\in\mathscr A_1\times\mathscr A_2$）时的特例。因此\cref{theorem:ch02-Tonelli}中 $\mu_1$ 与 $\mu_2$ 为 $\sigma$-有限的假设不能放松（见\cref{remark:ch02-sigma有限必要}）。
\end{remark}

对可积实值函数有下列结果。

\begin{theorem}[Fubini]\label{theorem:ch02-Fubini}
设 $(X_1,\mathscr A_1,\mu_1)$ 与 $(X_2,\mathscr A_2,\mu_2)$ 为 $\sigma$-有限测度空间。设 $f\in L^1(X_1\times X_2,\mathscr A_1\times\mathscr A_2,\mu_1\times\mu_2)$。则：
\begin{enumerate}[(i)]
\item 对 $\mu_1$-a.e. $x_1\in X_1$，$f_{x_1}\in L^1(X_2,\mathscr A_2,\mu_2)$；对 $\mu_2$-a.e. $x_2\in X_2$，$f_{x_2}\in L^1(X_1,\mathscr A_1,\mu_1)$；
\item 从 $X_1$ 到 $\mathbb R$ 的映射 $x_1\mapsto\int_{X_2}f_{x_1}\,\mathrm{d}\mu_2$ 与从 $X_2$ 到 $\mathbb R$ 的映射 $x_2\mapsto\int_{X_1}f_{x_2}\,\mathrm{d}\mu_1$ 分别属于 $L^1(X_1,\mathscr A_1,\mu_1)$ 与 $L^1(X_2,\mathscr A_2,\mu_2)$；
\item 有
\begin{align}
\int_{X_1\times X_2}f\,\mathrm{d}(\mu_1\times\mu_2)
=\int_{X_1}\mathrm{d}\mu_1\int_{X_2}f_{x_1}\,\mathrm{d}\mu_2
=\int_{X_2}\mathrm{d}\mu_2\int_{X_1}f_{x_2}\,\mathrm{d}\mu_1. \label{align:ch02-2.18}
\end{align}
\end{enumerate}
\end{theorem}

\cref{theorem:ch02-Fubini}中 $f$ 可积的假设不能放松，简单例子即表明。因此为应用 Fubini 定理，Tonelli 定理的如下推论是方便的。

\begin{theorem}\label{theorem:ch02-Fubini判别}
设 $(X_1,\mathscr A_1,\mu_1)$ 与 $(X_2,\mathscr A_2,\mu_2)$ 为 $\sigma$-有限测度空间。设 $f:X_1\times X_2\to\mathbb R$ 为 $\mathscr A_1\times\mathscr A_2$-可测。则下列命题等价：
\begin{enumerate}[(i)]
\item $f\in L^1(X_1\times X_2,\mathscr A_1\times\mathscr A_2,\mu_1\times\mu_2)$；
\item $\int_{X_1}\mathrm{d}\mu_1\int_{X_2}|f_{x_1}|\,\mathrm{d}\mu_2<\infty$；
\item $\int_{X_2}\mathrm{d}\mu_2\int_{X_1}|f_{x_2}|\,\mathrm{d}\mu_1<\infty$。
\end{enumerate}
\end{theorem}




\end{document}